\documentclass[12pt]{article}
\usepackage{amsmath, amsfonts, amssymb, amsthm, mathdots, mathrsfs, dsfont,systeme}
\usepackage[utf8]{inputenc}
\usepackage[spanish]{babel}
\usepackage{color}
\usepackage{fancyhdr}
\usepackage{enumerate}
\usepackage{hyperref}
\hypersetup{
    colorlinks=true,
    linkcolor=blue,
    filecolor=magenta,   
    citecolor=blue,
    urlcolor=cyan,
    pdftitle={Algebra lineal},
    pdfpagemode=FullScreen,
    }
\usepackage{amscd}
\usepackage{nomencl}
\usepackage{tikz}
\usepackage{tikz-cd}
\usepackage{amssymb}
\usepackage{blkarray}
\usepackage{multicol}
\usepackage[a4paper,margin=1cm,footskip=.5cm]{geometry}
\usepackage[all]{xy}
\usepackage{pst-node}
\usepackage{subfigure}
%\usepackage[table]{xcolor}
\usepackage{colortbl}
\newcommand{\cN}{\mathcal{N}}
\newcommand{\V}{\mathbb{V}}
\newcommand{\A}{\mathbb{A}}
\newcommand{\M}{\mathcal{M}}
\newcommand{\cH}{\mathcal{H}}
\newcommand{\F}{\mathcal{F}}
\newcommand{\G}{\mathcal{G}}
\newcommand{\cL}{\mathcal{L}}
\newcommand{\cO}{\mathcal{O}}
\newcommand{\Rn}{\mathbb{R}^n}
\DeclareMathOperator{\Hom}{Hom}
\DeclareMathOperator{\Ns}{NS}
\DeclareMathOperator{\im}{im}
\DeclareMathOperator{\Pic}{Pic}
\DeclareMathOperator{\Div}{Div}
\DeclareMathOperator{\Pdiv}{PDiv}
\DeclareMathOperator{\ddiv}{div}
\DeclareMathOperator{\Cl}{Cl}
\DeclareMathOperator{\CaDiv}{CaDiv}
\DeclareMathOperator{\Mor}{Mor}
\DeclareMathOperator{\Spec}{Spec}
\DeclareMathOperator{\id}{id}
\DeclareMathOperator{\Aut}{Aut}
\DeclareMathOperator{\End}{End}
\DeclareMathOperator{\CH}{CH}
\DeclareMathOperator{\IM}{Im}
\DeclareMathOperator{\Pf}{Pf}
\DeclareMathOperator{\Gl}{GL}
\DeclareMathOperator{\Sl}{SL}
\DeclareMathOperator{\St}{St}
\DeclareMathOperator{\tr}{Tr}
\DeclareMathOperator{\Gr}{Gr}
\DeclareMathOperator{\diag}{diag}
\DeclareMathOperator{\ord}{ord}
\DeclareMathOperator{\supp}{supp}
\DeclareMathOperator{\sign}{sign}
\DeclareMathOperator{\adj}{adj}
\DeclareMathOperator{\nul}{nul}
\DeclareMathOperator{\ran}{ran}
\newcommand{\calP}{\cal{P}}
\newcommand{\p}{\mathbb{P}}
\newcommand{\R}{\mathbb{R}}
\newcommand{\Z}{\mathbb{Z}}
\newcommand{\Q}{\mathbb{Q}}
\newcommand{\C}{\mathbb{C}}
\newcommand{\N}{\mathbb{N}}
\newcommand{\bbH}{\mathbb{H}}
\setlength{\parindent}{0pt}
\setlength{\textwidth}{6.5in}
\setlength{\textheight}{9.0in}
\setlength{\topmargin}{0in}
\setlength{\oddsidemargin}{0in}
\newtheoremstyle{exampstyle}
  {\topsep} % Space above
  {\topsep} % Space below
  {\itshape} % Body font
  {} % Indent amount
  {\bfseries} % Theorem head font
  {.} % Punctuation after theorem head
  {.5em} % Space after theorem head
  {} % Theorem head spec (can be left empty, meaning `normal')
\theoremstyle{exampstyle}
\newtheorem{thm}{Teorema}[section]
\newtheorem{prop}[thm]{Proposici\'on}
\newtheorem{lemma}[thm]{Lemma}
\newtheorem{cor}[thm]{Corolario}
\theoremstyle{definition}
\newtheorem{thme}{Expected Theorem}

\newtheorem{rec}[thm]{Recuerdo}
\newtheorem{conj}{Conjecture}
\newtheorem{defin}[thm]{Definici\'on}

\newtheoremstyle{break}%
{}{}%
{}{}%
{\bfseries}{}% % Note that final punctuation is omitted.
{\newline}{}
\theoremstyle{break}
\newtheorem{ejem}[thm]{Ejemplo}
\newtheorem{ejer}[thm]{Ejercicio}
\newtheorem{sol}[thm]{Soluci\'on}
\newtheorem{rmk}[thm]{Observaci\'on}
\numberwithin{equation}{section}
\newcommand{\gb}[1]{\operatorname{\textbf {#1}}}
\hypersetup{linkcolor=black}
\title{\'Algebra Lineal}
\author{Prof. Sebasti\'an Cea \\Prof. Camila Muñoz\\
}
\date{}
\pagestyle{fancy}

\lhead{\includegraphics[width = 4cm]{UFT.png}}
\setlength{\headsep}{0.6in}
\newcolumntype{a}{ >{\columncolor{yellow}} c }

\fancyhead[RO,LE]{
Prof: C. Muñoz.\\
Prof: A. Silva.\\
02/2023}}



\begin{document}
%\maketitle
%\newpage

\begin{minipage}[r]{15cm}


\end{minipage}

\begin{center}
\Large{\textbf{\'Algebra Lineal\\Tarea  }}
\end{center}

\section{Ejercicios}


\begin{enumerate}
\item Determine el o los valores del parámetro $\lambda$ para que el sistema no tenga solución, tenga una \'unica soluci\'on o tenga infinitas soluciones.



\begin{eqnarray*}
x + y +z &=& -2 \\
-\lambda x +3y + z &=& -7 \\
x+2 y + (\lambda +2)z &=& -5 \\
\end{eqnarray*}


    

%   Escribimos la matriz ampliada del sistema para luego comenzar a escalonar
% $$
%   \left(
% \begin{array}{ccc|c}
% k & 7 & 5  & 0 \\
% 1 & k & 1  & 0 \\
% 0 & 1 & 1  & -2 \\
% \end{array}
% \right) 
% \stackbin[  ]{\textcolor{blue}{f_{1}\longleftrightarrow f_{2}} }{\sim} 
%   \left(
% \begin{array}{ccc|c}
% 1 & k & 1  & 0 \\
% k & 7 & 5  & 0 \\
% 0 & 1 & 1  & -2 \\
% \end{array}
% \right) 
% \stackbin[  ]{\textcolor{blue}{f_{2}-kf_{1}} }{\sim} 
%   \left(
% \begin{array}{ccc|c}
% 1 & k & 1  & 0 \\
% 0 & 7-k^{2} & 5-k  & 0 \\
% 0 & 1 & 1  & -2 \\
% \end{array}
% \right) 
% $$

% $$
% \stackbin[  ]{\textcolor{blue}{f_{2}\longleftrightarrow f_{3}} }{\sim} 
%   \left(
% \begin{array}{ccc|c}
% 1 & k & 1  & 0 \\
% 0 & 1 & 1  & -2 \\
% 0 & 7-k^{2} & 5-k  & 0 \\
% \end{array}
% \right) 
% \stackbin[  ]{\textcolor{blue}{f_{3}-(7-k^{2})f_{2}} }{\sim} 
%   \left(
% \begin{array}{ccc|c}
% 1 & k & 1  & 0 \\
% 0 & 1 & 1  & -2 \\
% 0 & 0 & k^{2}-k-2  & 2(7-k^{2}) \\
% \end{array}
% \right)\textcolor{brown}{(1.4 puntos)} 
% $$
% Factorizando obtenemos
% $$
%   \left(
% \begin{array}{ccc|c}
% 1 & k & 1  & 0 \\
% 0 & 1 & 1  & -2 \\
% 0 & 0 & (k+1)(k-2)  & 2(\sqrt{7}-k)(\sqrt{7}+k) \\
% \end{array}
% \right) \textcolor{brown}{(0.8 puntos)}
% $$
%  Para que no tenga solución necesitamos que  $\rho(A)\neq \rho (A|B)$. Por lo tanto,  $k=-1$ o  $k=2$ \textcolor{brown}{(0,8 puntos)}, 

\item Obtenga los valores de $a$ y $b$ para que el vector  $u=(a+2b,a+3b,2)$ sea ortogonal a los vectores $v_1=(3,2,-1), v_2=(1,-2,-3) .$

 \item Es sabido que si $\alpha$ es valor propio de $A$ matriz asociada a la aplicación lineal, si y s\'olo si existe $v\neq 0$ tal que $Av=\alpha v$.
Si $H$ es una aplicación lineal invertible y $\alpha$ es valor propio de $H$, demuestre que $\dfrac{1}{\alpha}$ es valor propio de $H^{-1}$.


\item  Sea $T:\mathbb{R}^3\to \mathbb{R}^3$ tal que \\
$$T(1,0,0)=(1,-2,2),$$
$$T(0,1,0)=(2,-1,2),$$
$$T(0,0,1)=(0,0,1).$$
\begin{enumerate}
    \item Determine la matriz $[T]_{can}^B$, con $B=\{(1,0,0),(0,-1,1),(0,0,1)\}$.
\item Determine el polinomio caracter\'istico de la matriz anterior, sus valores y  vectores propios; subespacios propios de la matriz.
%\[A=\begin{pmatrix} 1 & 2  & 0  \\ 2 & 1 & 0\\ 0 & 1 & 1 \end{pmatrix}. \]
\item De ser posible, encuentre la matriz $P$, tal que $[T]_{can}^B=PDP^{-1}$, con $D$ matriz diagonal. En caso negativo, justifique.
\end{enumerate}



 
\end{enumerate}
\section*{Instrucciones}
\begin{itemize}
    \item Grupos de entre 4 y 5 personas.
    \item Entrega: d\'ia S\'abado 9 de diciembre a las 23:59.
\end{itemize}
\end{document}




